\section{近期研究（2026 年 8--11 月）}
\paragraph{Struct-GStream：用结构化高斯跟随基本运动，再补新出现区域}
2026.08~\cite{2608_01053}。Struct-GStream 将自由视点视频拆成可随时间更新的结构化高斯和用于补缺的自由高斯。前者沿近似局部刚体更新以复用已有结构，后者补足初始表示不足和之后出现的区域。原文比较跨帧整体补片与每帧独立生成，以区分新物体重建和时间连贯性。

论文在 N3DV、Meet Room 和 Google Immersive 给出对照，并报告所测百万像素场景平均每帧约 3.0 至 5.4 MB、在线训练每帧少于 10 秒、渲染超过 100 FPS。后两项分别是场景更新与显示，不是同一个端到端流延迟。作者指出首帧在稀疏角落和不可见室外区域重建不足会引起抖动，仍未达到实用流系统的全部要求。

\paragraph{D²-4DGS：把两类深度分别用于稀疏动态高斯的监督与补充}
2026.08~\cite{2608_01588}。少量同步相机只靠 RGB 拟合时，动态原语容易被放到只能解释训练视图的位置。D²-4DGS 对齐多视图与单目深度，先用验证过的深度做清理和约束，再将可靠多视图深度及部分单目估计用于提出缺失区域的增密候选。原有时序参数化保留，改动放在几何证据与 RGB-D 联合优化。

原文在 N3DV 六个、Technicolor 五个和 ENeRF-Outdoor 三个场景，分别让 2、3 或 4 路标定相机参与训练，其余有效相机留出。几何证据的益处在两路视图时较明显。结果针对同步标定相机和验证通过的深度，并不能把对齐的单目提议等同于实测深度。

\paragraph{FAST-GS：用傅里叶时间基表达动态高斯运动}
2026.08~\cite{2608_01958}。FAST-GS 以傅里叶分解替换时空高斯中的多项式运动参数，通过各轴的正弦和余弦系数表达时间变化，并加上运动正则以约束振荡和噪声。这是对动态高斯中心运动的表示调整，仍从多视图视频优化用于渲染的场景。

论文在 N3V 留出一台相机测视图泛化，在 Google Immersive 则用前 80\% 帧训练、后 20\% 帧测试。空间留出与时间留出的结果对应不同难度，需按协议分别解读。有限的频率基与有限观测时间也不能保证未见长时间行为的重建。

\paragraph{ASTRA：用二维轨迹帮助估计多相机时间偏移}
2026.08~\cite{2608_02006}。不同相机的帧率或采集起点不一致时，仅用颜色误差可把时间错位吸收到场景形状和运动中。ASTRA 在动态高斯上联合优化每台相机的时间偏移，用三维运动投影与观测轨迹对齐，并用动态和置信掩码减少不可靠匹配的影响。

作者将框架接入 4DGS 与 EDGS，在 Plenoptic 和 Blender 等数据评估 5 至 25 帧的偏移，并另用低纹理数据及其他实录验证。原文指出简单可靠性掩码在严重遮挡下可能失效，真实硬件的滚动快门、时钟漂移及曝光变化也超出当前常量偏移设置。

\paragraph{ACA-GS：让动态锚点按局部需求分配原语和特征}
2026.08~\cite{2608_04581}。锚点网络若为每个锚点生成相同数量的高斯，并分配相同特征容量，静态简单区域与变形复杂区域会消耗类似预算。ACA-GS 在 GIFStream 基础上允许锚点输出数量变化，把冗余高斯转给需要细节的区域，并通过学习到的通道掩码减少无用的锚点特征。

原文在 MPEG、Panoptic Sports 与 N3DV 比较渲染和表示大小，报告所测场景约 3.5 至 7.3 MB 的表示。三套协议的 LPIPS 骨干并不相同，不能跨表直接混算。容量选择和掩码强度也会改变压缩与画质的权衡。

\paragraph{RORA：把视频中的物体转成带关节的可模拟资产}
2026.08~\cite{2608_04842}。RORA 先从单段视频建立高斯和网格混合表示，再做几何分解。使用者参与把分解结果归组为可动部件、确认候选关节，高斯负责外观，网格提供可绑定部件及模拟几何。系统最后导出可交互资产，因此关节恢复是半自动流程。

作者在 PartNet-Mobility 对照 Articulate-Anything 和 ScrewSplat，并展示真实拍摄物体的模拟操作。对比中更新了 Articulate-Anything 的视觉语言模型，这一点影响复现实验定义。模拟演示也不能替代对不可见关节、柔性形变或网格误差下物理精度的验证。

\paragraph{ERF-GS：从不同相机位置的事件与 RGB 恢复快速运动}
2026.08~\cite{2608_08531}。多相机的事件和 RGB 输入不必从同一视点获得。ERF-GS 在动态 RGB 高斯骨干上引入两条事件线索，分别约束渲染及原语增密，使模糊或低帧率 RGB 缺少的运动信息可由另一位置的事件输入补充。这里的事件损失不要求同视点的清晰 RGB 真值。

论文在 Neu3D-dv 和 Nvidia-dv 比较，除整图指标还用 SAM2 标出运动区域计算动态 PSNR。作者说明主要多视图事件和模糊数据来自模拟，源 RGB 帧率限制合成事件的时间分辨率，真实高速运动和多传感器采集的差异仍需验证。

\paragraph{用黎曼流匹配生成和细化动态高斯参数}
2026.08~\cite{2608_18388}。题为 Depth Anything V4 的论文将单目视频特征送入黎曼流匹配模型，在 4D 高斯参数空间以常微分方程生成从粗到细的表示，再用渲染损失作测试时优化。该流程包含预训练先验与测试适配，部署成本不能等同于只作一次深度前馈。

原文用 Dynamic-Sim 和 KITTI Dynamic 分析动态重建，另以 ETH3D、ScanNet 做新视图测量。论文控制实验在相同预训练和测试细化条件下，相对确定性 MLP 报告 F 分数提高 0.044。预训练能否摊销取决于实际部署规模和测试优化开销，不能将实验设定直接外推。

\paragraph{按时间变化匹配四维高斯的空间滤波尺度}
2026.08~\cite{2608_21828}。动态高斯的运动会改变其跨相机的投影大小，使用固定抗混叠尺度可能在放大时显露空洞，在缩小时出现锯齿。该方法为原语估计时间与焦距深度比的联合分布，在给定时间据条件分布选取采样尺度，再对投影高斯施加平滑。

论文以 SARO-GS 为主要底座，在 Plenoptic Video 上测试低分辨率训练后的一到四倍输出，另用 D-NeRF 做缩小实验。原文指出为每个原语保存时间与焦距深度比的离散密度矩阵，内存随原语数和两个维度的分箱增长，大场景需要另计这项成本。

\paragraph{NemoSplat：在水下无标定视频中联合预测几何、瞬态与介质}
2026.08~\cite{2608_22888}。散射和游动物体会同时污染水下相机跟踪与静态地图。NemoSplat 从无标定视频经前馈网络联合估计相机、动态分离、可渲染高斯和介质参数，并以衰减和后向散射解释部分颜色变化，避免直接把这些变化写进静态几何。

作者用六条合成序列评价跟踪和渲染、另用 14 条真实序列检验新视图。合成平均 ATE 为 1.88 米，优于列出的两种前馈几何对照 VGGT 和 StreamVGGT；原文也指出极低可见度的 Deepsea 上光流 SLAM 仍能更稳定跟踪。真实渲染测试不能充当真实位姿误差，语义与介质模块另增加训练显存。

\paragraph{EvoGS：把逐帧形变改成带观测校正的时间递推}
2026.09~\cite{2609_00994}。EvoGS 为每个高斯保存持续形变状态，由历史校正值外推出当前运动，再用轻量网络从当前条件生成补充观测。校正权重取决于历史残差和运动变化，原语克隆与分裂也沿校正后的形变进行，以减少不确定原语继续增密。

论文将模块装到多种形变骨干，并在 NeRF-DS、HyperNeRF 的 vrig 子集和 Neu3D 比较。Neu3D 表中网格骨干 Grid4D 的平均 PSNR 为 30.75 dB，EvoGS 对应配置为 32.44 dB。短时外推假设采样时间间隔近似均匀，并要保留额外状态；这些条件限制了对非均匀长时序的直接推断。

\paragraph{UniFusion：把各相机各时刻的单目深度对齐后再训练动态高斯}
2026.09~\cite{2609_05888}。多个相机或时间的单目预测具有不同尺度与偏移，单帧深度不能直接当成统一的几何监督。UniFusion 用同步相机的稀疏 SfM 参照联合拟合时空深度场，通过低秩时间权重和沿射线预测的深度偏移维持跨帧共享，并将对齐结果用于动态高斯的初始化和优化监督。

论文在 Ego-Exo4D 按时间留出帧，在 EgoHuman 按相机留出视图，两个数据均使用已知相机位姿。深度对齐也有独立于高斯后端的重投影一致性实验。原文展示快速复杂运动仍可受变形后端限制，改善深度先验不等于取消运动模型的表达瓶颈。

\paragraph{FastFlowGS：用多来源对应与时序先验初始化每帧高斯}
2026.09~\cite{2609_16310}。流式动态重建若直接从上一帧复制位置，快速运动会给当前帧优化留下较大误差。FastFlowGS 从不同粒度的匹配、轨迹和光流提取带置信度的观测，将它们通过相机几何结合到三维位置，并以 Kalman 式时间先验给当前帧提供更好的初值，再继续图像优化。

原文在 CMU-Panoptic 六段数据分别用原帧率和五倍抽样测试，并在 Monaco4D 检验大位移室外场景。位置预测不更新旋转和尺度，所以快速旋转或形变仍要在本帧重新拟合；低纹理及沿深度方向运动也会削弱对应。不同指标和序列的改善不应汇成一个通用加速比。

\paragraph{以风驱振子检验动态高斯的物理外推}
2026.09~\cite{2609_17810}。作者将树的高斯按空间位置分为刚性部件，为每个部件学习受实测风速驱动的阻尼振子，再通过时间积分得到绕枢轴的转动。实验分别改变相机、时间与风力，区分同窗拟合、新视角表现和物理外推能力。

在三棵合成树上，学习形变的基线在未见相机上优于刚性振子；中等及密集树上，当时间或风力超出训练分布时，排序反过来。结果支持物理先验可帮助外推，同时作者未可靠恢复阻尼，模型没有分支耦合和部件内弯曲，并要求测试时同步提供风速。

\paragraph{SplashSplat：在同步多视角下分开重建飞溅液体与背景}
2026.09~\cite{2609_20818}。SplashSplat 用带标定和液体掩码的同步图像分别拟合动态液体高斯与静态背景，物理与几何结构只在有观测支撑的位置提供约束。最终渲染将两部分组合得到视频。

评估用同一七相机装置，五台训练、两台保留，测试视点位于训练视点之间，且按各相机固定裁剪比较。结果说明该协议下真实飞溅的新视角效果，不能直接外推自由相机与未观测物理轨迹。

\paragraph{4DGS-Fixer：用密集点初始化和视频先验细化四视图动态重建}
2026.09~\cite{2609_21176}。4DGS-Fixer 用多视图深度估计得到初始的稠密点云，再让视频修复先验迭代细化欠观测区域。视频结果为后续动态高斯优化提供额外约束，真实相机观测仍是检验生成结果的依据。

Neural3DV 测试沿用 4C4D 的协议，选取四个空间分散的训练相机，其余作测试。作者指出生成质量可能影响细化，并且迭代仍然耗费算力；结果不能解释为实时四图推断。

\paragraph{Dynamic Thermal Gaussians：用共享几何组织动态 RGB 和热红外}
2026.09~\cite{2609_24531}。该方法将 RGB 与热红外置于同一规范高斯几何，分别用模态属性、时变形变和结构感知路由渲染。模态专有的辐射变化可以由专门原语和属性解释，同时两类观测共享空间与相机参考。

论文在 DynamicRGBT-Scenes 的十二个真实序列中报告热红外 PSNR 平均 32.24 dB。输入依赖 RGB 视频获得的 COLMAP 位姿及与热帧的人工同步，评测的是配对动态新视角，不是以未标定的纯热红外完成重建。

\paragraph{ChronoFuseGS：融合不同时相并记录每个高斯的持续性}
2026.09~\cite{2609_31339}。ChronoFuseGS 将各时相分别训练的高斯场先配准到同一坐标系，再给每个高斯配置按时相变化的持续性属性。联合精调从各时相共享观测，并约束持续性数值的相邻时相变化，以免提高渲染分数的同时抹掉变化线索。渲染可以用颜色提示某个局部结构在哪些时相出现。

论文比较了独立训练和合并精调结果，整体图像指标的增幅较小，部分局部结构的改善更明显。参与者研究在二维预渲染图中检查变化标示。大范围变化可能使 SfM 对齐失败，单次采集中发生的结构或照明变化也不在当前建模范围内。

\paragraph{OneFixer：用同一次自回归生成轨迹训练街景视频修复}
2026.09~\cite{2609_32175}。OneFixer 将部署时的自回归历史纳入训练，在同一条生成轨迹上同时做流匹配与逐帧感知监督。它用单步生成修复输入的街景高斯渲染，训练目标兼顾帧质量和连续外观，而输出属于视频修复结果。

论文分别用 Waymo 和首尔驾驶数据训练，报告视频和逐帧指标，并说明细结构仍可能出现问题。不同数据域的模型训练分开进行，单步时延和图像一致性也不能直接解释为三维几何重建精度。

\paragraph{Endo-TSR：为内镜高斯的颜色和位移分别补充时间频谱}
2026.09~\cite{2609_32399}。Endo-TSR 在基础形变模型上为颜色和位移分别学习受限的傅里叶残差。颜色支路处理时间光度变化，位移支路由频谱先验约束，减少外观误差被挤进轨迹。多尺度拉普拉斯监督只作用在有效组织区域，以恢复细小纹理。

六段 EndoNeRF 和 StereoMIS 序列沿用既有七比一训练测试划分。定量结果评价渲染质量，消融中颜色建模带来的 PSNR 收益最大。这并不证明内部组织运动已准确恢复，且方法输入依赖组织掩码、深度先验和相机参数。

\paragraph{ProDyGS：以单目深度生成静止相机没有采到的代理视角}
2026.09~\cite{2609_32711}。静止或近乎静止的相机没有跨帧视差，难以用不同帧三角化动态表面。ProDyGS 先用单图模型预测每帧深度，把颜色随深度前向投影到相邻虚拟视角，再与真实相机观测一起约束规范高斯和时间形变网络。

在 DyNeRF，系统用单一相机序列训练，并在未见的真实相机视角评价。预测深度与 COLMAP 坐标的尺度平移对齐是该实验的评价设置。代理图像重投影的是同一帧已有观测，不能提供真实遮挡面证据。

\paragraph{DispFlow-GS：去掉相机运动后再监督形变位移}
2026.09~\cite{2609_36940}。DispFlow-GS 指出把渲染高斯流与二维光流直接拟合会混淆观测域。它用中间帧渲染分解相机与场景运动，在固定视点下用位移幅度约束形变。这样可以把图像质量指标与运动恢复指标分开检验。

论文在 NeRF-DS 观察到所用运动指标改善，但 PSNR 基本不变。HyperNeRF 存在因相机位姿误差而失败的场景。位移量监督也不等于知道逐像素对应，因此复杂运动中的精确轨迹仍难恢复。

\paragraph{EndoPrior-GS：用组织纹理引导稠密高斯生长与时间一致性}
2026.09~\cite{2609_37874}。EndoPrior-GS 以无需另训网络的联合纹理线索初始化内镜高斯，限制失真区域的冗余增密，并按邻近原语的局部纹理关系调整时序约束。在缺少可靠纹理的位置，方法无法获得同等强度的组织结构证据。

作者报告 EndoNeRF 和 SCARED 上的 Flow Error 下降。该指标来自 RAFT 的图像平面光流，并辅以渲染深度诊断，因为数据没有稠密三维组织运动真值。模糊、过曝、烟雾和失准掩码都可能削弱纹理先验。

\paragraph{用固定欧拉速度场和周期高斯输运重建水景}
2026.09~\cite{2609_38622}。该方法从具有相机视差的单目水景视频学习随位置变化、随时间固定的欧拉速度场。动态水体高斯沿固定速度场输运以产生可循环运动，余留的非周期外观变化另由时间残差描述。由掩码识别的静态高斯不随水体运动。

论文使用七段自采视频，将输入轨迹中的连续片段留作定量测试，并安排主观外观与运动比较。缺少同步多视角真值，因此留出片段的成绩不等于对独立角度的三维流速验证。

\paragraph{按表示机制理解动态重建的四维综述}
2026.09~\cite{2609_39960}。这篇综述以表示机制而非应用场景组织路线，分别梳理四维 NeRF 的隐式场、特征体和时间先验，以及高斯的显式时空原语、规范形变与逐帧训练。这个框架适合区分运动是直接写入基元，还是经共享规范坐标恢复。

性能章节汇总了动态场景数据、图像质量、几何和效率指标以及论文中的已报告结果。它能说明不同路线的取舍，但这些表格并非统一硬件、输入相机和训练协议下的重新实验。

\paragraph{以空间、时间和生成先验补全六相机动态场景}
2026.10~\cite{2610_04606}。该工作用前景掩码与背景深度约束不同区域，借助时间插值与光流约束投影运动，再在宽视角间隙内渲染虚拟视图。生成模型参考相邻已采图像修复虚拟渲染，结果作为伪观测循环回到四维高斯训练。

论文报告 SIGGRAPH Asia 2026 稀疏视图竞赛结果，并在验证集分别列出前景和整幅图像的消融分数。附加恢复可以弥补六相机覆盖的缺口，但虚拟视图不是实测，作者还将重建耗时与模型大小留作改进。

\paragraph{OuroWorld：把生成运动蒸馏为严格周期的四维高斯动画}
2026.10~\cite{2610_12461}。OuroWorld 从给定静态高斯场生成循环参考视频，再把跨视角生成的运动信号蒸馏到以傅里叶级数表示的周期形变。单独的漂移项吸收生成视图中不一致的部分，最终按周期形变驱动可渲染的三维基元。

作者在固定相机和绕行轨迹上分析运动活跃度、视觉质量、自然度和接缝，并辅以用户比较。这里没有真实动态真值。生成先验的错误可能进入最终场景，单参考视频在跨度大的新视角也更容易留下不一致。

\paragraph{Make It Flow：从静态高斯重建生成可循环流体动画}
2026.10~\cite{2610_13860}。Make It Flow 从输入参考图或静态高斯渲染图获得二维流动线索，再借高斯深度将不同相机的候选运动提升并融合成固定三维欧拉场。生成的局部形变作用到流体高斯上，通过对称时间处理衔接循环，不需要先采集动态多视图。

作者用七个随手采集的流动场景检查动画观感和成本，缺乏真实流场真值。反光水面或重建深度不一致会使提升的运动出错，静态模型里没表示的流体细节也不会因此成为实测几何。

\paragraph{Demiurge：自回归重渲染时持续合并四维点云记忆}
2026.10~\cite{2610_14685}。Demiurge 把已有单目视频和相机轨迹作为输入，从记忆中投影过去的几何及外观，为新的相机轨迹生成视频，再用补全深度把新观测写回四维点云。高斯重建可以使用这些生成多视图作为下游输入，核心记忆本身是点云。

论文分开评价生成视频的感知质量、跨轨迹同步和几何可重建性。依次执行外观生成、深度恢复与记忆更新，作者报告每条目标轨迹约需五分钟，当前不能视为交互帧率重渲染。

\paragraph{PUGS：将水下高斯位移建模为随场景变化的分布}
2026.10~\cite{2610_17936}。水下动态重建容易将折射及暂态视觉噪声误当作物体形变。PUGS 为每个高斯预测平均位移及方差，训练一条随机探索的渲染支路和一条确定性支路，并以自一致约束协调二者。学习的方差还用于优先补充较难可靠拟合的区域。

作者评价 UVEB、DRUVA 与 PhysGaia 共 22 个场景。强浑浊时图像误差和感知指标可能发生取舍，且不确定性退火仍需手工安排。这里的方差是重建中的可学习分量，不是经外部观测校准的运动误差概率。

\paragraph{SwitchSplat：将物体开关状态从连续时间中分离}
2026.10~\cite{2610_24812}。对可以打开、关闭或切换的物体，时间只是观测顺序，不能单独作为场景状态的索引。SwitchSplat 用共享静态底层储存不变部分，并为每个可交互物体维护可独立切换的二态高斯槽；可微槽表把各帧观测连到相应状态。这使新视图可以组合不同物体的状态，而不必重放某个采集时刻。

实验的常规十个场景与额外两组详细测试区分留出相机和新状态组合。方法当前假定状态为二元，采集约束在单人操作，槽位数量和大致交互窗口还由视觉语言模型给出初值。论文验证的是状态可控渲染，并未建立物理交互模型。

\paragraph{GTSM：在高斯空间支撑上附加连续热状态}
2026.10~\cite{2610_26918}。GTSM 将高斯支撑与其温度演变分开：三维或四维高斯后端描述原语位置和投影，附着在原语上的连续热状态由场景级潜在驱动和原语局部动态共同决定。对标定红外观测进行可微渲染，可以在已观测时间段查询新视角，也能让学习到的状态继续滚动。

实验分别使用真实 RHD 红外序列以及几何、温度共同变化的 D-NeRF 衍生数据，并区分段内重建与时外预测。几何运动会改变原语的空间覆盖，但几何变量不直接进入热演化方程，预测结果不等同于求解任意场景的热传导。

\paragraph{AEON：将视频生成潜空间接到显式四维高斯解码}
2026.10~\cite{2610_28631}。AEON 用一个四维重建模型从单目视频解码深度、高斯外观、动态和相机轨迹。为了让预训练视频扩散模型直接产生可查询的四维表示，另训练域适配器，把生成潜变量转换为重建潜变量。得到的显式高斯可沿指定视角和时间渲染。

论文分别检查动态重建、新轨迹渲染和生成视频，不同协议的重建与感知指标应分别解读。扩散模型序列长度限制当前生成时长，大幅转动相机可能暴露原视频从未生成的区域；跨窗口持续融合属于作者提出的扩展方向。

\paragraph{DUET4D：用当前高斯约束生成，再筛选伪观测回填}
2026.10~\cite{2610_35002}。DUET4D 先从当前四维高斯抽取跨时空对应，让冻结的单视图扩散模型恢复伪观测时沿这些锚点保持一致。生成结果再反向辅助建图，但只有置信度足够的差异才用于增密，与真实观测持续冲突的高斯则会被裁掉。

作者在三个稀疏视角动态数据上比较渲染效果。初始几何严重失真、运动太快或遮挡时间太长时，对应锚点的支持会下降。系统此时退回较独立的单视图恢复，可以限制错误传播，但无法保证重建真实不可见内容。

\paragraph{GaussianSpline：让每个动态高斯学习自己的时间分段}
2026.10~\cite{2610_35224}。GaussianSpline 为每个动态高斯分配分段三次贝塞尔轨迹，联合优化控制点与非均匀时间区间，因此固定段数可以向快运动处集中。接合点使用基于切向反射的几何连续性约束，在保持方向连续的同时允许相邻段速度变化。

论文在公开 N3DV、一个 Google Immersive 场景及两个合成序列评价。短时碰撞可能超出固定小段数的表达能力，而且各高斯轨迹独立学习，没有显式的跨原语刚体或关节约束。拓扑变化靠透明度和增密间接表示。

\paragraph{Topo-4DGS：在拓扑事件前后改变运动分支和表面支撑}
2026.11~\cite{2611_16012}。离散的拆开、合并、打开或闭合会改变物体部件运动及哪些表面可见。Topo-4DGS 按空间位置和事件进度对规范高斯分配分段运动，并在事件涉及的部位激活或停用相应高斯。新暴露的表面可用新原语补足，并绑定到所属部件的运动。

论文在 24 段第一人称视频上使用评审过的事件标注评价，不等同于全自动发现拓扑事件。当前适用对象是事件较局部且可近似分段刚性的部件。强形变、长时间手部遮挡和复合缓慢变化仍是难点。

\paragraph{ReCan-GS：用形变敏感度分离动态场景中的短时干扰}
2026.11~\cite{2611_16259}。ReCan-GS 同时维护可随时间形变的核心高斯和逐视图的干扰高斯。局部形变场的雅可比敏感度指出哪些运动难以由共享形变稳定解释，然后让这部分观测更倾向于进入干扰层，或者用于后处理裁除。

论文测试四个动态场景和四个合成烟雾场景。该分数并非经过校准的概率或语义标签，稳定出现的干扰可能保留，真实但复杂或观测偏少的运动也可能被削弱。作者还指出静态基准效果有限且缺乏专门烟雾移除基线。

\paragraph{OCLDA-GS：让流式高斯锚点沿物体边界传播运动}
2026.11~\cite{2611_23416}。OCLDA-GS 用固定锚点预算从已精调关键帧向后续视频帧传播运动。光流较强的物体组分得更多锚点；锚点插值限于组内，以减少物体边界穿插。每组刚性运动还作为局部非刚性偏移的条件，不要求每帧再拟合全部高斯。

论文在 N3DV、Google Immersive 等测试，域内对比仍有在线基线在 PSNR 上更强。组划分错误和缺少清楚主刚性变换的关节物体会影响传播，当前学习校正只更新位置，不纠正尺度、方向等其他属性。
